\begin{frame}{Teorema Mestre}
    \metroset{block=fill}

    \begin{Theorem}[Teorema Mestre]
        Sejam $a \geq 1$ e $b > 1$ inteiros constante, $f(n)$ uma função não negativa, e $T(n)$ definida nos inteiros não-negativos pela recorrência
            \[ T(n) = aT(n/b) + f(n) \ , \]
        interpretando $n/b$ como $\lceil n/b \rceil$ ou $\lfloor n/b \rfloor$. Então $T(n)$ tem os seguintes limites:
            \begin{itemize} \pause
                \item[(i)] Se $f(n) = O(n^{\log_b a - \varepsilon})$ para alguma constante $\varepsilon > 0$, então $T(n) = \Theta(n^{\log_b a})$. \pause
                
                \item[(ii)] Se $f(n) = \Theta(n^{\log_b a})$, então $T(n) = \Theta(n^{\log_b a} \log n)$. \pause

                \item[(iii)] Se $f(n) = \Omega(n^{\log_b a + \varepsilon})$ para alguma constante $\varepsilon > 0$, e se $af(n/b) \leq c f(n)$ para alguma constante $c < 1$ e todo $n$ suficientemente grande, então $T(n) = \Theta(f(n))$.
            \end{itemize}
    \end{Theorem}

    
\end{frame}